\documentclass[12pt,dvipdfmx]{jsarticle}

% A0 portrait poster template for pLaTeX + dvipdfmx.
\usepackage[
  paperwidth=841mm,
  paperheight=1189mm,
  top=35mm,
  bottom=35mm,
  left=35mm,
  right=35mm
]{geometry}

\usepackage{amsmath,amssymb,bm}
\usepackage{graphicx}
\usepackage{booktabs}
\usepackage{multicol}
\usepackage{tikz}
\usepackage{xcolor}
\usepackage{tcolorbox}
\tcbuselibrary{skins,breakable}
\usepackage{url}
\usepackage{hyperref}

\definecolor{posterblue}{RGB}{20,62,112}
\definecolor{postergreen}{RGB}{31,116,90}
\definecolor{postergray}{RGB}{245,247,250}
\definecolor{posterline}{RGB}{210,218,226}

\pagestyle{empty}
\setlength{\parindent}{0pt}
\setlength{\columnsep}{22mm}
\setlength{\columnseprule}{0pt}

\renewcommand{\baselinestretch}{1.12}
\newcommand{\postertitle}[1]{%
  {\fontsize{54}{64}\selectfont\bfseries #1\par}
}
\newcommand{\posterauthor}[1]{%
  \vspace{8mm}{\fontsize{28}{36}\selectfont #1\par}
}
\newcommand{\posteraffiliation}[1]{%
  \vspace{3mm}{\fontsize{22}{30}\selectfont #1\par}
}

\newtcolorbox{posterblock}[1]{
  enhanced,
  breakable,
  colback=white,
  colframe=posterline,
  coltitle=white,
  colbacktitle=posterblue,
  title={#1},
  fonttitle=\fontsize{25}{31}\selectfont\bfseries,
  fontupper=\fontsize{20}{30}\selectfont,
  boxrule=1.2pt,
  arc=2mm,
  left=7mm,
  right=7mm,
  top=6mm,
  bottom=7mm,
  before skip=0mm,
  after skip=12mm
}

\begin{document}

\begin{tikzpicture}[remember picture,overlay]
  \fill[postergray] (current page.south west) rectangle (current page.north east);
  \fill[posterblue] (current page.north west) rectangle ([yshift=-72mm]current page.north east);
\end{tikzpicture}

\vspace*{-8mm}
\begin{minipage}{0.74\textwidth}
  \color{white}
  \postertitle{ポスタータイトル}
  \posterauthor{著者名}
  \posteraffiliation{所属・研究科・メールアドレス}
\end{minipage}
\hfill
\begin{minipage}{0.20\textwidth}
  \raggedleft
  \color{white}
  {\fontsize{22}{28}\selectfont 学会名・日付など}
\end{minipage}

\vspace{34mm}
\color{black}

\begin{multicols}{3}

\begin{posterblock}{研究背景}
ここに研究背景を書く。問題設定、先行研究との差分、なぜ重要かを短くまとめる。

\vspace{4mm}
\begin{itemize}
  \item 背景となる制度・市場・モデル。
  \item 既存研究で分かっていること。
  \item 本研究が扱う未解決の点。
\end{itemize}
\end{posterblock}

\begin{posterblock}{目的}
本研究の目的を一文から数文で明確に書く。

\[
  \text{Research Question:}\quad
  \text{ここに主要な問いを置く}
\]
\end{posterblock}

\begin{posterblock}{モデル}
記号、仮定、制約を整理する。数式が多い場合は、ポスターでは最小限に絞る。

\[
  f : \mathcal{P} \to \mathcal{A}, \qquad
  \mu = f(P)
\]

必要なら図や表で流れを示す。
\end{posterblock}

\begin{posterblock}{手法}
証明、シミュレーション、実証分析など、使った方法を書く。

\begin{enumerate}
  \item データ・入力を準備する。
  \item メカニズムまたは推定手法を適用する。
  \item 評価指標で結果を比較する。
\end{enumerate}
\end{posterblock}

\begin{posterblock}{主要結果}
もっとも重要な結論を大きく、少ない数で示す。

\begin{center}
\begin{tabular}{lcc}
\toprule
条件 & 方法 A & 方法 B \\
\midrule
ケース 1 & 0.82 & 0.91 \\
ケース 2 & 0.76 & 0.88 \\
\bottomrule
\end{tabular}
\end{center}

表の下に、読み取れる含意を一、二文で書く。
\end{posterblock}

\begin{posterblock}{図を入れる場所}
\begin{center}
\fbox{%
  \begin{minipage}[c][105mm][c]{0.82\linewidth}
    \centering
    \fontsize{22}{28}\selectfont
    \texttt{\textbackslash includegraphics} で図を挿入
  \end{minipage}
}
\end{center}

例:
\begin{center}
\texttt{\textbackslash includegraphics[width=\textbackslash linewidth]\{figures/result.pdf\}}
\end{center}
\end{posterblock}

\begin{posterblock}{考察}
結果の解釈、政策的含意、既存研究との違いを書く。ポスターでは文章を長くしすぎず、箇条書きを中心にする。
\end{posterblock}

\begin{posterblock}{結論}
\begin{itemize}
  \item 結論 1。
  \item 結論 2。
  \item 今後の課題。
\end{itemize}
\end{posterblock}

\begin{posterblock}{参考文献}
必要に応じて \texttt{references.bib} を使う。

\vspace{3mm}
例: 著者名 (年), 論文タイトル, 雑誌名。
\end{posterblock}

\end{multicols}

\end{document}
